% Lessons 50--62: arrays, linked lists, bits, matrices, trees, and heaps.

\patternintoc{Prefix and Suffix Products}
\problemstart{Product of Array Except Self}{238}{Medium}{Prefix and suffix products}{Array}

\subsection{The Problem}
Return an array where position \(i\) contains the product of every input value
except \texttt{nums[i]}. Do not use division.

\subsection{Combine Products from Both Sides}
First store the product strictly to the left of each index. A reverse scan
keeps the product strictly to the right and multiplies it into the answer.

\begin{invariantbox}{the current answer combines all known outside values}
After the forward pass, \texttt{answer[i]} is the left prefix product. During
the reverse pass, \texttt{suffix} is the product right of \(i\), so their
product excludes exactly the current value.
\end{invariantbox}

\begin{visualbox}{Left and right products meet at one index}
\centering
\begin{tikzpicture}[font=\small,>={Stealth[length=2.2mm,width=1.55mm]}]
  \foreach \x/\v in {0/1,1/2,2/3,3/4}{\node[arraycell] (a\x) at (1.2*\x,0) {\v};}
  \node[draw=orange,very thick,fit=(a2),inner sep=1mm] {};
  \node[above=5mm of a2,text=orange,font=\bfseries] {leave out index 2};
  \node[below=9mm of a0,xshift=-2mm,draw=teal,fill=softteal,rounded corners,
    minimum width=27mm] (left) {left product: $1\times2=2$};
  \node[below=9mm of a3,xshift=12mm,draw=orange,fill=softorange,rounded corners,
    minimum width=23mm] (right) {right product: $4$};
  \node[below=11mm of left,xshift=22mm,draw=navy,fill=softblue,rounded corners]
    {combine: $2\times4=8$};
\end{tikzpicture}
\end{visualbox}

\subsection{Pseudocode}
\begin{lstlisting}[style=pseudocode]
prefix = 1
for i from left to right:
    answer[i] = prefix
    prefix *= nums[i]
suffix = 1
for i from right to left:
    answer[i] *= suffix
    suffix *= nums[i]
return answer
\end{lstlisting}

\subsection{C++ Implementation}
\begin{lstlisting}[style=compactcode,caption={Complete solution: Product of Array Except Self}]
#include <vector>

using namespace std;

class Solution {
public:
    vector<int> productExceptSelf(vector<int>& nums) {
        vector<int> answer(nums.size(), 1);
        int prefix = 1;
        for (int i = 0; i < static_cast<int>(nums.size()); ++i) {
            answer[i] = prefix; prefix *= nums[i];
        }
        int suffix = 1;
        for (int i = static_cast<int>(nums.size()) - 1; i >= 0; --i) {
            answer[i] *= suffix; suffix *= nums[i];
        }
        return answer;
    }
};
\end{lstlisting}

\subsection{Time, Space, and Edge Cases}
Time is \(O(n)\); excluding output, space is \(O(1)\). Zero values require no
special case: the products naturally become zero in the appropriate positions.

\lessonexpansion
  {For \texttt{[1,2,3,4]}, the left pass writes \([1,1,2,6]\). The right
   pass multiplies by suffixes \(24,12,4,1\), producing
   \([24,12,8,6]\). Neither pass includes the value at its own index.}
  {Before the right pass, \texttt{answer[i]} equals the product strictly left
   of \(i\). Multiplying by the running product strictly right of \(i\) gives
   exactly the required product of every other position.}
  {Check one zero, two zeroes, negative values, two elements, and products
   equal to one. Use wider arithmetic if adapting beyond the stated bounds.}
  {I first store every left-side product in the output. A reverse pass
   multiplies each entry by its right-side product, achieving linear time
   without division or extra arrays.}

\chaptercheckpoint{Each output excludes exactly one array position.}
  {A stored left product and a rolling right product.}
  {Multiplying the current value before writing its prefix.}
  {$O(n)$ time and $O(1)$ extra space.}


\patternintoc{Linked-List Pointers}
\problemstart{Remove Nth Node From End of List}{19}{Medium}{Fixed pointer gap}{Linked list}

\subsection{The Problem}
Remove the \(n\)-th node from the end of a singly linked list and return the
head.

\subsection{Create an n-Node Gap}
A dummy node handles deletion of the original head. Advance \texttt{fast}
\(n+1\) steps from the dummy, then move both pointers until \texttt{fast}
reaches null. \texttt{slow} then precedes the node to delete.

\begin{invariantbox}{fast stays n+1 links ahead of slow}
The fixed gap makes the remaining distance after \texttt{fast} reaches null
equal to the required offset from the end. So \texttt{slow->next} is
the target node.
\end{invariantbox}

\begin{visualbox}{The pointer gap locates the previous}
\centering
\begin{tikzpicture}[font=\small,>={Stealth[length=2.2mm,width=1.55mm]}]
  \node[dsnode] (d) {dummy}; \node[dsnode,right=7mm of d] (a) {1};
  \node[dsnode,right=7mm of a] (b) {2}; \node[dsnode,active,right=7mm of b] (c) {3};
  \node[dsnode,right=7mm of c] (e) {4}; \node[dsnode,right=7mm of e] (f) {5};
  \draw[flowarrow] (d)--(a)--(b)--(c)--(e)--(f);
  \node[below=7mm of b,text=teal]{slow}; \node[above=7mm of f,text=orange]{fast};
  \node[below=7mm of c,text=orange]{remove for \(n=3\)};
\end{tikzpicture}
\end{visualbox}

\subsection{Pseudocode}
\begin{lstlisting}[style=pseudocode]
dummy.next = head; fast = slow = dummy
move fast forward n + 1 steps
while fast exists:
    move fast and slow one step
slow.next = slow.next.next
return dummy.next
\end{lstlisting}

\subsection{C++ Implementation}
\begin{lstlisting}[style=compactcode,caption={Complete solution: Remove Nth Node From End}]
using namespace std;

class Solution {
public:
    ListNode* removeNthFromEnd(ListNode* head, int n) {
        ListNode dummy(0); dummy.next = head;
        ListNode* fast = &dummy; ListNode* slow = &dummy;
        for (int i = 0; i <= n; ++i) fast = fast->next;
        while (fast != nullptr) { fast = fast->next; slow = slow->next; }
        slow->next = slow->next->next;
        return dummy.next;
    }
};
\end{lstlisting}

\subsection{Time and Space}
One pass uses \(O(L)\) time and \(O(1)\) space. The dummy node makes removing
the first real node identical to every other deletion.

\lessonexpansion
  {For \texttt{1->2->3->4->5} with \(n=2\), advance \texttt{fast} three
   links beyond the dummy. Move both pointers until \texttt{fast} becomes
   null; \texttt{slow} then precedes node 4, which is bypassed.}
  {The initial gap is \(n+1\) links. Moving both pointers keeps it, so
   when \texttt{fast} reaches null, \texttt{slow.next} is exactly the
   \(n\)-th node from the end.}
  {Check a one-node list, removing the head, removing the tail, \(n=1\), and
   \(n\) equal to the list length.}
  {A dummy node eliminates the head-deletion special case. I create a fixed
   pointer gap, advance both pointers together, and bypass the node after the
   slow pointer.}

\chaptercheckpoint{A list position is measured from the tail.}
  {Two pointers separated by $n+1$ links.}
  {Using an $n$-link gap and landing on the target instead of its previous.}
  {$O(L)$ time and $O(1)$ space.}


\problemstart{Reorder List}{143}{Medium}{Split, reverse, and weave}{Linked list}

\subsection{The Problem}
Transform \(L_0\to L_1\to\cdots\to L_n\) into
\(L_0\to L_n\to L_1\to L_{n-1}\to\cdots\) in place.

\subsection{Three Pointer Phases}
Use slow and fast pointers to find the middle, reverse the second half, then
alternate nodes from the first half and the reversed second half.

\begin{invariantbox}{the woven prefix already has final order}
During merging, every node before the two current pointers is fixed in the
required outside-in order. Saved next pointers keep both unmerged suffixes.
\end{invariantbox}

\begin{visualbox}{Reverse the back half, then alternate}
\centering
\begin{tikzpicture}[font=\small,>={Stealth[length=2.2mm,width=1.55mm]}]
  \node[dsnode] (a) {1}; \node[dsnode,right=7mm of a] (b) {2};
  \node[dsnode,right=7mm of b] (c) {3}; \node[dsnode,right=7mm of c] (d) {4};
  \node[dsnode,right=7mm of d] (e) {5};
  \draw[flowarrow] (a)--(b)--(c)--(d)--(e);
  \node[below=12mm of c,dsnode,active] (out) {1 $\to$ 5 $\to$ 2 $\to$ 4 $\to$ 3};
  \draw[flowarrow,teal] (c)--(out);
\end{tikzpicture}
\end{visualbox}

\subsection{Pseudocode}
\begin{lstlisting}[style=pseudocode]
find the middle with slow and fast pointers
reverse the second half
first = head; second = reversed half
while second exists:
    save both next pointers
    link first -> second -> next first
    advance both halves
\end{lstlisting}

\subsection{C++ Implementation}
\begin{lstlisting}[style=compactcode,caption={Complete solution: Reorder List}]
using namespace std;

class Solution {
public:
    void reorderList(ListNode* head) {
        if (head == nullptr || head->next == nullptr) return;
        ListNode *slow = head, *fast = head;
        while (fast->next != nullptr && fast->next->next != nullptr) {
            slow = slow->next; fast = fast->next->next;
        }
        ListNode* second = slow->next; slow->next = nullptr;
        ListNode* previous = nullptr;
        while (second != nullptr) {
            ListNode* next = second->next; second->next = previous;
            previous = second; second = next;
        }
        ListNode* first = head; second = previous;
        while (second != nullptr) {
            ListNode* nextFirst = first->next;
            ListNode* nextSecond = second->next;
            first->next = second; second->next = nextFirst;
            first = nextFirst; second = nextSecond;
        }
    }
};
\end{lstlisting}

\subsection{Time and Space}
Each phase is linear, so total time is \(O(n)\) and space is \(O(1)\). Cutting
the first half before reversing prevents an accidental cycle.

\lessonexpansion
  {For \texttt{1->2->3->4->5}, slow stops at 3. Cut after 3 and reverse
   \texttt{4->5} into \texttt{5->4}. Weave 1 with 5, then 2 with 4; the first
   half's remaining node 3 stays last, producing
   \texttt{1->5->2->4->3}.}
  {The middle split creates two independent lists. Reversal places the original
   tail at the front of the second half. Each weaving step appends the next
   unused front node followed by the next unused back node, exactly matching
   the required order.}
  {Check two nodes, three nodes, even and odd lengths, repeated values, and a
   one-node list. Verify that the final node points to null and no cycle is
   introduced.}
  {I solve the problem in three linear phases: find the middle, reverse the
   second half, and alternate the two halves. Saved next pointers prevent the
   unmerged suffixes from being lost.}

\chaptercheckpoint{List nodes must be woven from alternating ends.}
  {Middle pointers, a reversed second half, and two merge cursors.}
  {Failing to disconnect the halves before weaving.}{$O(n)$ time and $O(1)$ space.}


\problemstart{Reverse Linked List}{206}{Easy}{Iterative pointer reversal}{Linked list}

\subsection{The Problem}
Reverse every link in a singly linked list and return the new head.

\subsection{Keep Next Before Rewiring}
Keep a reversed prefix and an untouched suffix. Save the current node's
next pointer, redirect it toward the prefix, and advance both boundaries.

\begin{invariantbox}{previous heads the fully reversed prefix}
Before each iteration, \texttt{previous} is a correctly reversed list of all
processed nodes, while \texttt{current} heads the untouched remainder.
\end{invariantbox}

\begin{visualbox}{One link changes direction per step}
\centering
\begin{tikzpicture}[font=\small,>={Stealth[length=2.2mm,width=1.55mm]}]
  \node[dsnode,visited] (a) {1}; \node[dsnode,active,right=12mm of a] (b) {2};
  \node[dsnode,right=12mm of b] (c) {3};
  \draw[flowarrow,teal] (b)--(a); \draw[flowarrow] (b)--(c);
  \node[above=6mm of a]{reversed prefix}; \node[above=6mm of c]{saved next};
\end{tikzpicture}
\end{visualbox}

\subsection{Pseudocode}
\begin{lstlisting}[style=pseudocode]
previous = null; current = head
while current exists:
    nextNode = current.next
    current.next = previous
    previous = current
    current = nextNode
return previous
\end{lstlisting}

\subsection{C++ Implementation}
\begin{lstlisting}[style=compactcode,caption={Complete solution: Reverse Linked List}]
using namespace std;

class Solution {
public:
    ListNode* reverseList(ListNode* head) {
        ListNode* previous = nullptr;
        while (head != nullptr) {
            ListNode* next = head->next;
            head->next = previous;
            previous = head;
            head = next;
        }
        return previous;
    }
};
\end{lstlisting}

\subsection{Time and Space}
Time is \(O(n)\), space is \(O(1)\). Rewiring before saving the next pointer
loses the unprocessed suffix.

\lessonexpansion
  {For \texttt{1->2->3}, save node 2 before changing node 1 to point to null.
   Then make 2 point to 1 and 3 point to 2. The final \texttt{previous}
   pointer is the new head 3.}
  {Before each iteration, \texttt{previous} heads the reversed prefix and
   \texttt{head} heads the untouched suffix. Saving the suffix link before
   rewiring keeps all nodes and grows the reversed prefix by one.}
  {Check a null list, one node, two nodes, repeated values, and a long list
   where pointer loss would be visible.}
  {I keep a reversed prefix and an unprocessed suffix. Each iteration
   saves the next node, reverses one pointer, and advances both boundaries.}

\chaptercheckpoint{Every next pointer must reverse direction.}
  {A reversed prefix, current node, and saved next node.}
  {Overwriting the only pointer to the remaining list.}{$O(n)$ time and $O(1)$ space.}


\patternintoc{Bit Manipulation}
\problemstart{Reverse Bits}{190}{Easy}{Fixed-width bit scan}{Unsigned integer}

\subsection{The Problem}
Reverse the order of all 32 bits in an unsigned integer.

\begin{invariantbox}{the result contains the reversed processed suffix}
After \(i\) iterations, the result's lowest \(i\) construction steps contain
the first \(i\) input bits in reverse order. Shifting before insertion opens
the next output position.
\end{invariantbox}

\begin{visualbox}{Bits leave the right and enter the left-to-right result}
\centering
\begin{tikzpicture}[font=\ttfamily\small,>={Stealth[length=2.2mm,width=1.55mm]}]
  \node[dsnode] (in) {input: 101100};
  \node[dsnode,active,right=25mm of in] (bit) {take 0};
  \node[dsnode,visited,right=25mm of bit] (out) {result: ...0};
  \draw[flowarrow] (in)--(bit)--(out);
\end{tikzpicture}
\end{visualbox}

\subsection{Pseudocode}
\begin{lstlisting}[style=pseudocode]
answer = 0
repeat 32 times:
    answer = (answer << 1) OR (n AND 1)
    n = n >> 1
return answer
\end{lstlisting}

\subsection{C++ Implementation}
\begin{lstlisting}[style=compactcode,caption={Complete solution: Reverse Bits}]
#include <cstdint>

using namespace std;

class Solution {
public:
    uint32_t reverseBits(uint32_t n) {
        uint32_t result = 0;
        for (int i = 0; i < 32; ++i) {
            result = (result << 1) | (n & 1U);
            n >>= 1;
        }
        return result;
    }
};
\end{lstlisting}

\subsection{Time and Space}
Exactly 32 iterations give \(O(1)\) time and space. Processing all positions,
including leading zeros, is essential.

\lessonexpansion
  {Using an 8-bit illustration, \texttt{00010110} is consumed from the right.
   Appending its bits to the result produces \texttt{01101000}. The real
   implementation repeats the same operation exactly 32 times.}
  {At iteration \(i\), the result contains the lowest \(i\) input bits in
   reversed order. Shifting before appending makes room for the next bit;
   exactly 32 rounds also place original leading zeroes correctly.}
  {Check zero, one, only the highest bit, all bits set, and alternating bit
   patterns.}
  {I read one low bit at a time, shift the result left, append that bit, and
   shift the input right. A fixed 32 iterations keep leading zeroes.}

\chaptercheckpoint{Reverse a fixed-width bit representation.}
  {An input shift register and an output running value.}
  {Stopping when the input becomes zero and losing leading positions.}{32 steps and $O(1)$ space.}


\patternintoc{Matrix Transformations}
\problemstart{Rotate Image}{48}{Medium}{Transpose and reverse rows}{Square matrix}

\subsection{The Problem}
Rotate an \(n\times n\) matrix 90 degrees clockwise in place.

\subsection{Two In-place Reflections}
Transposing swaps \((r,c)\) with \((c,r)\). Reversing every row then moves
the transposed columns into their clockwise positions.

\begin{invariantbox}{each swap fixes a symmetric pair}
The transpose loop visits only cells above the diagonal, so every off-diagonal
pair is exchanged exactly once. Row reversal then completes the mapping
\((r,c)\mapsto(c,n-1-r)\).
\end{invariantbox}

\begin{visualbox}{Transpose, then reverse each row}
\centering
\begin{tikzpicture}[font=\small]
  \matrix[matrix of nodes,nodes={arraycell,minimum width=7mm,minimum height=7mm},
    row sep=-\pgflinewidth,column sep=-\pgflinewidth] (a) {1&2&3\\4&5&6\\7&8&9\\};
  \matrix[right=24mm of a,matrix of nodes,nodes={arraycell,minimum width=7mm,minimum height=7mm},
    row sep=-\pgflinewidth,column sep=-\pgflinewidth] (b) {1&4&7\\2&5&8\\3&6&9\\};
  \matrix[right=24mm of b,matrix of nodes,nodes={arraycell,minimum width=7mm,minimum height=7mm},
    row sep=-\pgflinewidth,column sep=-\pgflinewidth] (c) {7&4&1\\8&5&2\\9&6&3\\};
  \draw[flowarrow] (a.east) -- node[edgelabel,above=2mm,fill=white]{transpose} (b.west);
  \draw[flowarrow,teal] (b.east) -- node[edgelabel,above=2mm,fill=white]{reverse rows} (c.west);
  \node[below=5mm of a]{start}; \node[below=5mm of c,text=teal,font=\bfseries]{rotated};
\end{tikzpicture}
\end{visualbox}

\subsection{Pseudocode}
\begin{lstlisting}[style=pseudocode]
for every pair above the main diagonal:
    swap matrix[row][col] with matrix[col][row]
for every row:
    reverse the row
\end{lstlisting}

\subsection{C++ Implementation}
\begin{lstlisting}[style=compactcode,caption={Complete solution: Rotate Image}]
#include <algorithm>
#include <vector>

using namespace std;

class Solution {
public:
    void rotate(vector<vector<int>>& matrix) {
        int n = static_cast<int>(matrix.size());
        for (int r = 0; r < n; ++r)
            for (int c = r + 1; c < n; ++c)
                swap(matrix[r][c], matrix[c][r]);
        for (auto& row : matrix) reverse(row.begin(), row.end());
    }
};
\end{lstlisting}

\subsection{Time and Space}
Every cell is touched a constant number of times: \(O(n^2)\) time and
\(O(1)\) extra space.

\lessonexpansion
  {For \(\begin{smallmatrix}1&2&3\\4&5&6\\7&8&9\end{smallmatrix}\),
   transposition gives
   \(\begin{smallmatrix}1&4&7\\2&5&8\\3&6&9\end{smallmatrix}\).
   Reversing each row then yields the clockwise rotation
   \(\begin{smallmatrix}7&4&1\\8&5&2\\9&6&3\end{smallmatrix}\).}
  {Transposition maps \((r,c)\) to \((c,r)\). Reversing each row then maps it
   to \((c,n-1-r)\), which is exactly the coordinate transformation for a
   90-degree clockwise rotation.}
  {Check \(1\times1\), \(2\times2\), odd and even dimensions, negative
   entries, and avoid swapping diagonal pairs twice.}
  {I decompose the rotation into two in-place symmetries: transpose across the
   main diagonal, then reverse every row.}

\chaptercheckpoint{A square matrix must rotate in place.}
  {Symmetric transpose pairs followed by row boundaries.}
  {Swapping both halves of the transpose and undoing the work.}{$O(n^2)$ time.}


\patternintoc{Tree Comparison and Serialization}
\problemstart{Same Tree}{100}{Easy}{Paired recursion}{Two binary trees}

\subsection{The Problem}
Return whether two binary trees have identical structure and equal values at
matching nodes.

\begin{invariantbox}{each call compares matching subtrees}
Two null nodes match; exactly one null node fails; otherwise equal roots and
matching left and right subtree pairs are needed and enough.
\end{invariantbox}

\begin{visualbox}{Matching nodes are compared in pairs}
\centering
\begin{tikzpicture}[font=\small]
  \node[trienode] (a) {1}; \node[trienode,below left=10mm and 8mm of a] (a1) {2};
  \node[trienode,below right=10mm and 8mm of a] (a2) {3};
  \node[trienode,right=40mm of a] (b) {1}; \node[trienode,below left=10mm and 8mm of b] (b1) {2};
  \node[trienode,below right=10mm and 8mm of b] (b2) {3};
  \draw[-,navy] (a)--(a1) (a)--(a2) (b)--(b1) (b)--(b2);
  \node[draw=teal,rounded corners,fill=softteal,below=18mm of a,xshift=20mm,
    align=center] {compare $(1,1)$, then $(2,2)$ and $(3,3)$\\all values and positions match};
\end{tikzpicture}
\end{visualbox}

\subsection{Pseudocode}
\begin{lstlisting}[style=pseudocode]
same(p, q):
    if both are null: return true
    if exactly one is null or values differ: return false
    return same(p.left, q.left) AND same(p.right, q.right)
\end{lstlisting}

\subsection{C++ Implementation}
\begin{lstlisting}[style=compactcode,caption={Complete solution: Same Tree}]
using namespace std;

class Solution {
public:
    bool isSameTree(TreeNode* p, TreeNode* q) {
        if (p == nullptr || q == nullptr) return p == q;
        return p->val == q->val && isSameTree(p->left, q->left) &&
               isSameTree(p->right, q->right);
    }
};
\end{lstlisting}

\subsection{Time and Space}
At most all nodes are compared: \(O(n)\) time and \(O(h)\) recursion space.

\lessonexpansion
  {Compare roots first. If both contain 1, recursively compare their left
   children and then their right children. A missing child on only one side
   fails immediately even if every remaining value matches.}
  {Two rooted trees are equal exactly when their roots match and their
   matching left and right subtrees are equal. The base cases handle
   matching nulls and structural mismatches, proving the recursive test.}
  {Check two null trees, one null tree, equal values with different shapes,
   one differing leaf, and identical skewed trees.}
  {I compare matching nodes. Two nulls match; one null or unequal values
   fail; otherwise both matching child pairs must also match.}

\chaptercheckpoint{Two trees must match in both shape and values.}
  {Pairs of matching subtree roots.}
  {Comparing traversal values while ignoring null structure.}{$O(n)$ time and $O(h)$ space.}


\problemstart{Serialize and Deserialize Binary Tree}{297}{Hard}{Preorder with null markers}{Tree and string stream}

\subsection{The Problem}
Convert a binary tree into a string and reconstruct the identical tree from
that string.

\subsection{Keep Shape with Null Tokens}
Preorder values alone are ambiguous. Emit a marker for every null child. The
decoder consumes tokens in the same recursive order, so each value receives
the exact left and right structure it originally had.

\begin{invariantbox}{each recursive decode consumes one complete subtree}
The current token is the root description. A null marker consumes one token;
a value consumes itself followed by exactly the tokens for its left and right
subtrees.
\end{invariantbox}

\begin{visualbox}{Null markers make tree shape explicit}
\centering
\begin{tikzpicture}[font=\small,>={Stealth[length=2.2mm,width=1.55mm]}]
  \node[trienode] (r) {1}; \node[trienode,below left=10mm and 12mm of r] (a) {2};
  \node[trienode,below right=10mm and 12mm of r] (b) {3};
  \draw[flowarrow] (r)--(a); \draw[flowarrow] (r)--(b);
  \node[dsnode,right=28mm of r] {\texttt{1,2,\#,\#,3,\#,\#}};
\end{tikzpicture}
\end{visualbox}

\subsection{Pseudocode}
\begin{lstlisting}[style=pseudocode]
serialize(node):
    if null: append null marker
    else append value, serialize left, serialize right
deserialize(tokens):
    read next token
    if null marker: return null
    node = new node(token)
    node.left = deserialize(tokens)
    node.right = deserialize(tokens)
    return node
\end{lstlisting}

\subsection{C++ Implementation}
\begin{lstlisting}[style=compactcode,caption={Complete solution: Serialize and Deserialize Tree}]
#include <sstream>
#include <string>

using namespace std;

class Codec {
    void write(TreeNode* node, ostringstream& out) {
        if (node == nullptr) { out << "# "; return; }
        out << node->val << ' '; write(node->left, out); write(node->right, out);
    }
    TreeNode* read(istringstream& in) {
        string token; in >> token;
        if (token == "#") return nullptr;
        TreeNode* node = new TreeNode(stoi(token));
        node->left = read(in); node->right = read(in); return node;
    }
public:
    string serialize(TreeNode* root) { ostringstream out; write(root, out); return out.str(); }
    TreeNode* deserialize(string data) { istringstream in(data); return read(in); }
};
\end{lstlisting}

\subsection{Time and Space}
Both operations visit every node and null child: \(O(n)\) time, \(O(n)\)
output/tree space, and \(O(h)\) recursion space.

\lessonexpansion
  {For a root 1 with left child 2 and right child 3, preorder serialization
   emits \texttt{1 2 \# \# 3 \# \#}. During decoding, token 1 creates the
   root; the next complete token group constructs its left subtree, and the
   remaining group constructs its right subtree.}
  {Preorder plus null markers is clear. A value token determines one
   node and recursively consumes exactly two subtree encodings; a null marker
   ends one child. Working upward from the leaves shows that
   decoding recreates the original shape and values.}
  {Check an empty tree, one node, negative values, a left-only chain, a
   right-only chain, repeated values, and a tree whose shape would be ambiguous
   without null markers.}
  {I serialize in preorder and write an explicit token for every null pointer.
   The decoder reads the same grammar recursively, so each call consumes one
   complete subtree. Both operations are linear.}

\chaptercheckpoint{A tree must round-trip through a string without losing shape.}
  {Preorder tokens including explicit null markers.}
  {Serializing only values and making shapes ambiguous.}{$O(n)$ time per operation.}


\patternintoc{Matrix State}
\problemstart{Set Matrix Zeroes}{73}{Medium}{First-row and first-column markers}{Matrix}

\subsection{The Problem}
If a matrix cell is zero, set its entire row and column to zero in place.

\subsection{Store Markers Inside the Matrix}
Use the first cell of each row and column as a zero marker, plus one Boolean
for whether the original first column contained zero. Apply markers after the
discovery pass so newly written zeros do not spread incorrectly.

\begin{invariantbox}{markers describe original zero rows and columns}
After discovery, \texttt{matrix[r][0]} and \texttt{matrix[0][c]} are zero
exactly when row \(r\) or column \(c\) must be cleared.
\end{invariantbox}

\begin{visualbox}{A zero writes one row marker and one column marker}
\centering
\begin{tikzpicture}[font=\small]
  \matrix[matrix of nodes,nodes={arraycell},row sep=-\pgflinewidth,column sep=-\pgflinewidth] (a) {
    1&1&1\\1&0&1\\1&1&1\\};
  \node[draw=orange,very thick,fit=(a-2-2),inner sep=1mm] {};
  \matrix[right=22mm of a,matrix of nodes,nodes={arraycell},row sep=-\pgflinewidth,column sep=-\pgflinewidth] (b) {
    1&0&1\\0&0&0\\1&0&1\\};
  \draw[flowarrow] (a.east)--node[edgelabel,above,fill=white]{mark row and column}(b.west);
  \node[below=5mm of a]{original};
  \node[below=5mm of b,text=teal,font=\bfseries]{final matrix};
\end{tikzpicture}
\end{visualbox}

\subsection{Pseudocode}
\begin{lstlisting}[style=pseudocode]
record whether first row or first column originally contains zero
use first row and first column to mark zeroed rows and columns
zero all marked interior cells
zero the first row and first column when originally required
\end{lstlisting}

\subsection{C++ Implementation}
\begin{lstlisting}[style=compactcode,caption={Complete solution: Set Matrix Zeroes}]
#include <vector>

using namespace std;

class Solution {
public:
    void setZeroes(vector<vector<int>>& a) {
        int rows = a.size(), cols = a[0].size(); bool firstCol = false;
        for (int r = 0; r < rows; ++r) {
            if (a[r][0] == 0) firstCol = true;
            for (int c = 1; c < cols; ++c)
                if (a[r][c] == 0) a[r][0] = a[0][c] = 0;
        }
        for (int r = rows - 1; r >= 0; --r) {
            for (int c = cols - 1; c >= 1; --c)
                if (a[r][0] == 0 || a[0][c] == 0) a[r][c] = 0;
            if (firstCol) a[r][0] = 0;
        }
    }
};
\end{lstlisting}

\subsection{Time and Space}
Time is \(O(mn)\), extra space \(O(1)\). The separate first-column flag
resolves the shared marker at \texttt{matrix[0][0]}.

\lessonexpansion
  {For \texttt{[[1,1,1],[1,0,1],[1,1,1]]}, discovery writes zero into the
   first cell of row 1 and column 1. The second pass clears every interior cell
   whose row or column marker is zero, producing
   \texttt{[[1,0,1],[0,0,0],[1,0,1]]}.}
  {Markers are written only while reading the original matrix. Once discovery
   is complete, a cell must become zero exactly when its row marker or column
   marker is zero. The separate first-column flag resolves the dual meaning of
   the top-left marker.}
  {Check one cell, one row, one column, a zero at \((0,0)\), zeros only in the
   first row or first column, multiple zeroes, and an all-zero matrix.}
  {I reuse the first row and first column as marker storage, which avoids
   allocating separate Boolean arrays. I delay mutation of the remaining cells
   until all original zero locations have been recorded.}

\chaptercheckpoint{Zeros propagate across complete rows and columns.}
  {First-row/column markers and one first-column flag.}
  {Clearing while still discovering original zeros.}{$O(mn)$ time and $O(1)$ space.}


\problemstart{Spiral Matrix}{54}{Medium}{Shrinking boundaries}{Matrix}

\subsection{The Problem}
Return matrix values in clockwise spiral order.

\begin{invariantbox}{the boundaries enclose every unvisited cell}
At the start of each round, all cells outside \texttt{top}, \texttt{bottom},
\texttt{left}, and \texttt{right} have been emitted exactly once. Traversing
the perimeter and shrinking its sides keeps this condition.
\end{invariantbox}

\begin{visualbox}{Four boundaries peel one rectangular layer}
\centering
\begin{tikzpicture}[font=\small,>={Stealth[length=2.2mm,width=1.55mm]}]
  \matrix[matrix of nodes,nodes={arraycell},row sep=-\pgflinewidth,column sep=-\pgflinewidth] (m) {
    1&2&3\\8&9&4\\7&6&5\\};
  \foreach \cell/\num in {m-1-1/1,m-1-2/2,m-1-3/3,m-2-3/4,
    m-3-3/5,m-3-2/6,m-3-1/7,m-2-1/8,m-2-2/9}
    \node[font=\scriptsize,text=orange,anchor=north east] at (\cell.north east) {\num};
  \draw[orange,very thick,rounded corners=2mm]
    ([xshift=-2mm,yshift=2mm]m-1-1.north west) rectangle
    ([xshift=2mm,yshift=-2mm]m-3-3.south east);
  \node[right=10mm of m-2-3,align=left]
    {read the outside: $1,2,3,4,5,6,7,8$\\then read the center: $9$};
\end{tikzpicture}
\end{visualbox}

\subsection{Pseudocode}
\begin{lstlisting}[style=pseudocode]
set top, bottom, left, right boundaries
while boundaries are valid:
    read top left-to-right; increment top
    read right top-to-bottom; decrement right
    if rows remain: read bottom right-to-left; decrement bottom
    if columns remain: read left bottom-to-top; increment left
return output
\end{lstlisting}

\subsection{C++ Implementation}
\begin{lstlisting}[style=compactcode,caption={Complete solution: Spiral Matrix}]
#include <vector>

using namespace std;

class Solution {
public:
    vector<int> spiralOrder(vector<vector<int>>& a) {
        vector<int> out; int top=0, bottom=a.size()-1, left=0, right=a[0].size()-1;
        while (top <= bottom && left <= right) {
            for (int c=left;c<=right;++c) out.push_back(a[top][c]); ++top;
            for (int r=top;r<=bottom;++r) out.push_back(a[r][right]); --right;
            if (top <= bottom) { for (int c=right;c>=left;--c) out.push_back(a[bottom][c]); --bottom; }
            if (left <= right) { for (int r=bottom;r>=top;--r) out.push_back(a[r][left]); ++left; }
        }
        return out;
    }
};
\end{lstlisting}

\subsection{Time and Space}
Every cell is emitted once: \(O(mn)\) time and \(O(1)\) extra space.
The checks before the bottom and left traversals prevent duplicates in a
single remaining row or column.

\lessonexpansion
  {For a \(3\times4\) matrix, read the top row, right column, bottom row in
   reverse, and left column upward. The boundaries then describe only the
   inner row, which is read left-to-right without revisiting a cell.}
  {After each side is emitted, moving its boundary inward removes exactly
   those cells from the remaining rectangle. Conditional bottom and left
   passes prevent duplication when that rectangle collapses to one dimension.}
  {Check one cell, one row, one column, \(2\times n\), \(m\times2\), and
   non-square matrices.}
  {I keep four boundaries around the unvisited rectangle and consume its
   sides clockwise. Each cell leaves the rectangle exactly once.}

\chaptercheckpoint{Read a matrix by successive outer layers.}
  {Four shrinking boundaries.}
  {Traversing the final row or column twice.}{$O(mn)$ time.}


\patternintoc{Tree Containment}
\problemstart{Subtree of Another Tree}{572}{Easy}{Tree matching at every root}{Binary trees}

\subsection{The Problem}
Return whether \texttt{subRoot} occurs as an identical subtree within
\texttt{root}.

\subsection{Separate Search from Equality}
At every main-tree node, test full tree equality. If it fails, recursively
search the left and right subtrees.

\begin{invariantbox}{every possible matching root is examined}
The search call returns true exactly when a matching root occurs in its
current subtree. The equality helper verifies both values and null structure.
\end{invariantbox}

\begin{visualbox}{Choose a possible root, then compare its full shape}
\centering
\begin{tikzpicture}[font=\small]
  \node[trienode] (r) {3}; \node[trienode,active,below left=11mm and 15mm of r] (a) {4};
  \node[trienode,below right=11mm and 15mm of r] (b) {5};
  \node[trienode,visited,below left=10mm and 7mm of a] (c) {1};
  \node[trienode,visited,below right=10mm and 7mm of a] (d) {2};
  \draw[-,navy] (r)--(a) (r)--(b) (a)--(c) (a)--(d);
  \node[draw=teal,dashed,fit=(a)(c)(d),inner sep=3mm]{ };
\end{tikzpicture}
\end{visualbox}

\subsection{Pseudocode}
\begin{lstlisting}[style=pseudocode]
isSubtree(root, subRoot):
    if subRoot is null: return true
    if root is null: return false
    if sameTree(root, subRoot): return true
    return isSubtree(root.left, subRoot)
           OR isSubtree(root.right, subRoot)
\end{lstlisting}

\subsection{C++ Implementation}
\begin{lstlisting}[style=compactcode,caption={Complete solution: Subtree of Another Tree}]
using namespace std;

class Solution {
    bool same(TreeNode* a, TreeNode* b) {
        if (a == nullptr || b == nullptr) return a == b;
        return a->val == b->val && same(a->left,b->left) && same(a->right,b->right);
    }
public:
    bool isSubtree(TreeNode* root, TreeNode* subRoot) {
        if (root == nullptr) return false;
        return same(root, subRoot) || isSubtree(root->left, subRoot) ||
               isSubtree(root->right, subRoot);
    }
};
\end{lstlisting}

\subsection{Time and Space}
Worst-case time is \(O(mn)\) when equality is retried at many similar nodes;
recursion uses \(O(h)\) combined depth. Hashing subtrees can improve repeated
comparison but is more complex.

\lessonexpansion
  {If \texttt{subRoot} begins with value 4, visit every value-4 node in the
   larger tree. At each choice, compare the complete left and right shapes;
   one extra descendant is enough to reject that choice.}
  {Any valid subtree must begin at some node of the larger tree. The outer
   recursion examines every possible root, and the equality helper accepts
   exactly when the entire rooted structure and all values match.}
  {Check identical trees, a one-node subtree, repeated choice values,
   matching values with different shape, and a possible subtree larger than the host.}
  {I treat every node as a possible subtree root and run an exact tree-equality
   check there. If it fails, I continue recursively on both child subtrees.}

\chaptercheckpoint{Find one tree as an exact rooted region of another.}
  {A choice-root search plus a paired equality test.}
  {Matching only values while ignoring extra descendants.}{$O(mn)$ worst-case time.}


\patternintoc{Bitwise Arithmetic}
\problemstart{Sum of Two Integers}{371}{Medium}{XOR and carry}{Bits}

\subsection{The Problem}
Add two integers without using \texttt{+} or \texttt{-}.

\subsection{Separate Partial Sum and Carry}
XOR adds bits without carry. AND finds positions generating carry, shifted one
place left. Repeat until no carry remains.

\begin{invariantbox}{partial sum plus carry equals the original total}
At every iteration, \texttt{a} contains the carry-free partial sum and
\texttt{b} contains all pending carries. Redistributing them with XOR and AND
keeps their mathematical sum.
\end{invariantbox}

\begin{visualbox}{XOR writes sum bits; AND creates carries}
\centering
\begin{tikzpicture}[font=\ttfamily\small,>={Stealth[length=2.2mm,width=1.55mm]}]
  \node[dsnode] (a) {0101}; \node[dsnode,below=7mm of a] (b) {0011};
  \node[dsnode,right=20mm of a] (x) {XOR: 0110};
  \node[dsnode,active,right=20mm of b] (c) {carry: 0010};
  \draw[flowarrow] (a)--(x); \draw[flowarrow] (b)--(c);
\end{tikzpicture}
\end{visualbox}

\subsection{Pseudocode}
\begin{lstlisting}[style=pseudocode]
while b is not zero:
    carry = (a AND b) << 1
    a = a XOR b
    b = carry
return a
\end{lstlisting}

\subsection{C++ Implementation}
\begin{lstlisting}[style=compactcode,caption={Complete solution: Sum of Two Integers}]
#include <cstdint>

using namespace std;

class Solution {
public:
    int getSum(int a, int b) {
        uint32_t x = static_cast<uint32_t>(a);
        uint32_t y = static_cast<uint32_t>(b);
        while (y != 0) { uint32_t carry = (x & y) << 1; x ^= y; y = carry; }
        return static_cast<int>(x);
    }
};
\end{lstlisting}

\subsection{Time and Space}
For 32-bit integers, at most 32 carry movements give \(O(1)\) time and space.
Unsigned intermediates make left shifting well-defined.

\lessonexpansion
  {For 5 and 3, XOR gives 6, the sum without carries. AND identifies the
   shared low bit and shifting produces carry 2. Repeating with 6 and 2 gives
   partial sum 4 and carry 4, then final result 8.}
  {XOR adds each bit without carry, while shifted AND contains exactly the
   carries required at the next position. Repeating until no carry remains is
   is equivalent to ordinary binary addition.}
  {Check zero, two positive values, opposite signs, two negative values, and
   results near the 32-bit limits.}
  {I separate binary addition into a carry-free sum using XOR and a carry mask
   using shifted AND. I repeat until the carry mask becomes zero.}

\chaptercheckpoint{Perform addition through bit operations.}
  {A carry-free XOR sum and shifted carry bits.}
  {Shifting negative signed integers.}{At most 32 iterations.}


\patternintoc{Frequencies and Heaps}
\problemstart{Top K Frequent Elements}{347}{Medium}{Frequency heap}{Hash map and min heap}

\subsection{The Problem}
Return the \(k\) values occurring most often in an array.

\subsection{Count, Then Keep Only k Choices}
Build a frequency map. A min heap ordered by frequency stores at most \(k\)
entries; when it grows larger, remove the least frequent choice.

\begin{invariantbox}{the heap contains the best k frequencies seen so far}
After each map entry is processed, any discarded element has frequency no
greater than every retained element. So the final heap contains exactly a
valid top-\(k\) set.
\end{invariantbox}

\begin{visualbox}{The smallest retained frequency is easiest to evict}
\centering
\begin{tikzpicture}[font=\small]
  \node[dsnode] (map) {value:count\\\(1:3,\ 2:2,\ 3:1\)};
  \node[trienode,active,right=30mm of map] (h) {\(2:2\)};
  \node[trienode,below=11mm of h] (a) {\(1:3\)};
  \draw[-,thick,navy] (h)--(a);
  \node[above=5mm of h,font=\bfseries,align=center]
    {min heap for \(k=2\)\\root has the smallest retained count};
  \node[right=24mm of h,dsnode,active,align=center]
    {\(3:1\)\\count 1 was removed};
  \node[below=4mm of a,text=teal] {retained counts: 2 and 3};
\end{tikzpicture}
\end{visualbox}

\subsection{Pseudocode}
\begin{lstlisting}[style=pseudocode]
count every value in a hash map
minHeap = empty
for each (value, frequency):
    push pair into heap
    if heap size exceeds k: pop smallest frequency
return the values remaining in the heap
\end{lstlisting}

\subsection{C++ Implementation}
\begin{lstlisting}[style=compactcode,caption={Complete solution: Top K Frequent Elements}]
#include <queue>
#include <unordered_map>
#include <utility>
#include <vector>

using namespace std;

class Solution {
public:
    vector<int> topKFrequent(vector<int>& nums, int k) {
        unordered_map<int,int> count; for (int x:nums) ++count[x];
        using Entry=pair<int,int>;
        priority_queue<Entry,vector<Entry>,greater<Entry>> heap;
        for (const auto& entry : count) {
            heap.push({entry.second, entry.first});
            if ((int)heap.size() > k) heap.pop();
        }
        vector<int> answer; while(!heap.empty()){answer.push_back(heap.top().second);heap.pop();}
        return answer;
    }
};
\end{lstlisting}

\subsection{Time and Space}
For \(n\) values and \(u\) distinct values, time is \(O(n+u\log k)\), with
\(O(u+k)\) space. Any output order is accepted.

\lessonexpansion
  {For \texttt{[1,1,1,2,2,3]} with \(k=2\), the map stores
   \(1\mapsto3\), \(2\mapsto2\), and \(3\mapsto1\). A size-two min heap
   eventually retains frequencies three and two, so the returned values are
   1 and 2.}
  {After processing any number of map entries, evicting the smallest frequency
   whenever the heap exceeds \(k\) leaves the best \(k\) entries seen so far.
   Induction over the entries proves the final heap is a valid top-\(k\) set.}
  {Check \(k=1\), \(k\) equal to the number of distinct values, negative
   values, tied frequencies, one input value, and heavily duplicated input.}
  {I count first, then keep only \(k\) choices in a min heap. Keeping
   the weakest retained choice at the top makes replacement efficient and
   avoids sorting all distinct values.}

\chaptercheckpoint{Select values by frequency rather than numeric magnitude.}
  {A count map and a size-$k$ min heap.}
  {Building a max heap of all values when only $k$ are needed.}{$O(n+u\log k)$ time.}